\subsection{\texorpdfstring{\(\mathcal{L}(\mathbf{E};\mathbf{F})\) 与 \(\mathcal{L}(\hat{\mathbf{E}};\mathbf{F})\) 之间的关系}{\textbackslash mathcal\{L\}(\textbackslash mathbf\{E\};\textbackslash mathbf\{F\}) 与 \textbackslash mathcal\{L\}(\textbackslash hat\{\textbackslash mathbf\{E\}\};\textbackslash mathbf\{F\}) 之间的关系}}

设 E 和 F 是两个 Hausdorff 局部凸空间，并假设 F 完备；设 \(\hat{E}\) 是 E 的完备化。由于从 E 到 F 中的每个连续线性映射 u 都唯一地扩张为一个连续线性映射 \(\overline{u}\)（从 \(\hat{E}\) 到 F 中），我们可以把空间 \(\mathcal{L}(E;F)\) 与 \(\mathcal{L}(\hat{E};F)\) 通过映射 \(u\mapsto\overline{u}\) 等同起来。此外，设 S 是 E 的有界子集的一个族；\(\mathcal{L}(E;F)\) 上的 S-拓扑与 \(\mathcal{L}(\hat{E};F)\) 上的 S-拓扑重合，也与 \(\hat{S}\) -拓扑重合，其中 \(\hat{S}\) 表示 S 中诸集在 \(\hat{E}\) 中的闭包所成的族。

例如，若 E 是赋范的，则 \(\mathcal{L}(\mathrm{E};\mathrm{F})\) 上的有界收敛拓扑与 \(\mathcal{L}(\hat{\mathrm{E}};\mathrm{F})\) 上的有界收敛拓扑相同：因为 \(\hat{E}\) 的每个有界子集都包含在 E 的某个有界子集的闭包中。由于 \(\hat{E}\) 的单位球是 E 的单位球的闭包，由公式 (3)（III, 第 14 页）可得：若 F 是 Banach 空间，则映射 \(u\mapsto\overline{u}\) 是从 \(\mathcal{L}(\mathrm{E};\mathrm{F})\) 到 \(\mathcal{L}(\hat{\mathrm{E}};\mathrm{F})\) 上的一个等距。

我们注意到，若 E 不是赋范空间，则可能存在 \(\hat{E}\) 的有界子集，它不含于 E 的任何有界子集的闭包之中（例如，当 E 是无限维 Banach 空间的弱对偶时）；然而，若 E 可度量化且满足第一可数公理，则此事成立（III, 第 39 页, 习题 16）。

